\documentclass[a4paper, 14pt]{extarticle}

% Поля
%--------------------------------------
\usepackage{geometry}
\geometry{a4paper,tmargin=2cm,bmargin=2cm,lmargin=3cm,rmargin=1cm}
%--------------------------------------

\usepackage{float}
%Russian-specific packages
%--------------------------------------
\usepackage[T2A]{fontenc}
\usepackage[utf8]{inputenc} 
\usepackage[english, main=russian]{babel}
%--------------------------------------

\usepackage{textcomp}

% Красная строка
%--------------------------------------
\usepackage{indentfirst}               
%--------------------------------------             
\usepackage[most]{tcolorbox}
\tcbuselibrary{listings,breakable}
\newtcblisting{breakablelisting}{
  listing only,
  breakable,
  enhanced,
  frame hidden,
  colback=white,
  listing options={
    mathescape=false,
    texcl=false,
    inputencoding=utf8,
    literate={\$}{{\$}}1
  }
}


%Graphics
%--------------------------------------
\usepackage{graphicx}
\usepackage{wrapfig}
%--------------------------------------
\graphicspath{{./lab5/}}
% Полуторный интервал
%--------------------------------------
\linespread{1.3}                    
%--------------------------------------

%Выравнивание и переносы
%--------------------------------------
% Избавляемся от переполнений
\sloppy
% Запрещаем разрыв страницы после первой строки абзаца
\clubpenalty=10000
% Запрещаем разрыв страницы после последней строки абзаца
\widowpenalty=10000
%--------------------------------------

%Списки
\usepackage{enumitem}

%Подписи
\usepackage{caption} 

%Гиперссылки
\usepackage{hyperref}

\hypersetup {
	unicode=true
}

%Рисунки
%--------------------------------------
\DeclareCaptionLabelSeparator*{emdash}{~--- }
\captionsetup[figure]{labelsep=emdash,font=onehalfspacing,position=bottom}
%--------------------------------------

\usepackage{tempora}

%Листинги
%--------------------------------------
\usepackage{listings}
\lstset{
  basicstyle=\ttfamily\footnotesize, 
  %basicstyle=\footnotesize\AnkaCoder,        % the size of the fonts that are used for the code
  breakatwhitespace=false,         % sets if automatic breaks shoulbd only happen at whitespace
  breaklines=true,                 % sets automatic line breaking
  captionpos=t,                    % sets the caption-position to bottom
  inputencoding=utf8,
  frame=single,                    % adds a frame around the code
  keepspaces=true,                 % keeps spaces in text, useful for keeping indentation of code (possibly needs columns=flexible)
  keywordstyle=\bf,       % keyword style
  numbers=left,                    % where to put the line-numbers; possible values are (none, left, right)
  numbersep=5pt,                   % how far the line-numbers are from the code
  xleftmargin=25pt,
  xrightmargin=25pt,
  showspaces=false,                % show spaces everywhere adding particular underscores; it overrides 'showstringspaces'
  showstringspaces=false,          % underline spaces within strings only
  showtabs=false,                  % show tabs within strings adding particular underscores
  stepnumber=1,                    % the step between two line-numbers. If it's 1, each line will be numbered
  tabsize=2,                       % sets default tabsize to 8 spaces
  title=\lstname                   % show the filename of files included with \lstinputlisting; also try caption instead of title
}
%--------------------------------------

%%% Математические пакеты %%%
%--------------------------------------
\usepackage{amsthm,amsfonts,amsmath,amssymb,amscd}  % Математические дополнения от AMS
\usepackage{mathtools}                              % Добавляет окружение multlined
\usepackage[perpage]{footmisc}
%--------------------------------------

%--------------------------------------
%			НАЧАЛО ДОКУМЕНТА
%--------------------------------------

\begin{document}

%--------------------------------------
%			ТИТУЛЬНЫЙ ЛИСТ
%--------------------------------------
\begin{titlepage}
\thispagestyle{empty}
\newpage


%Шапка титульного листа
%--------------------------------------
\vspace*{-60pt}
\hspace{-65pt}
\begin{minipage}{0.3\textwidth}
\hspace*{-20pt}\centering
\includegraphics[width=\textwidth]{emblem}
\end{minipage}
\begin{minipage}{0.67\textwidth}\small \textbf{
\vspace*{-0.7ex}
\hspace*{-6pt}\centerline{Министерство науки и высшего образования Российской Федерации}
\vspace*{-0.7ex}
\centerline{Федеральное государственное автономное образовательное учреждение }
\vspace*{-0.7ex}
\centerline{высшего образования}
\vspace*{-0.7ex}
\centerline{<<Московский государственный технический университет}
\vspace*{-0.7ex}
\centerline{имени Н.Э. Баумана}
\vspace*{-0.7ex}
\centerline{(национальный исследовательский университет)>>}
\vspace*{-0.7ex}
\centerline{(МГТУ им. Н.Э. Баумана)}}
\end{minipage}
%--------------------------------------

%Полосы
%--------------------------------------
\vspace{-25pt}
\hspace{-35pt}\rule{\textwidth}{2.3pt}

\vspace*{-20.3pt}
\hspace{-35pt}\rule{\textwidth}{0.4pt}
%--------------------------------------

\vspace{1.5ex}
\hspace{-35pt} \noindent \small ФАКУЛЬТЕТ\hspace{80pt} <<Информатика и системы управления>>

\vspace*{-16pt}
\hspace{47pt}\rule{0.83\textwidth}{0.4pt}

\vspace{0.5ex}
\hspace{-35pt} \noindent \small КАФЕДРА\hspace{50pt} <<Теоретическая информатика и компьютерные технологии>>

\vspace*{-16pt}
\hspace{30pt}\rule{0.866\textwidth}{0.4pt}
  
\vspace{11em}

\begin{center}
\Large {\bf Лабораторная работа № 8} \\ 
\large {\bf по курсу <<Методы оптимизации>>} \\
\large <<Метод Пауэлла>> 
\end{center}\normalsize

\vspace{8em}


\begin{flushright}
  {Студент группы ИУ9-82Б Шемякин В.А. \hspace*{15pt}\\ 
  \vspace{2ex}
  Преподаватель Посевин Д. П.\hspace*{15pt}}
\end{flushright}

\bigskip

\vfill
 

\begin{center}
\textsl{Москва 2026}
\end{center}
\end{titlepage}
%--------------------------------------
%		КОНЕЦ ТИТУЛЬНОГО ЛИСТА
%--------------------------------------

\renewcommand{\ttdefault}{pcr}

\setlength{\tabcolsep}{3pt}
\newpage
\setcounter{page}{2}

\section{Задание}\label{Sect::task}
Необходимо реализовать метод Пауэлла
\section{Результаты}\label{Sect::res}

Исходный код программы представлен в листинге 1.

\begin{breakablelisting}
using Plots
using LinearAlgebra

rosenbrock(x) = (1.0 - x[1])^2 + 100.0 * (x[2] - x[1]^2)^2
rastrigin(x) = 20.0 + (x[1]^2 - 10.0 * cos(2π * x[1])) + (x[2]^2 - 10.0 * cos(2π * x[2]))
schwefel(x) = 418.9829 * 2.0 - x[1] * sin(sqrt(abs(x[1]))) - x[2] * sin(sqrt(abs(x[2])))

function find_interval(f, x, d; h=0.1, factor=2.0, max_iter=20)
    a, fa = 0.0, f(x)
    b, fb = h, f(x + h * d)
    if fb > fa
        h = -h
        b, fb = h, f(x + h * d)
        fb > fa && return (-abs(h), abs(h))
    end
    c = b + h
    for _ in 1:max_iter
        fc = f(x + c * d)
        fc > fb && return (min(a, c), max(a, c))
        a, fa, b, fb = b, fb, c, fc
        h *= factor
        c = b + h
    end
    return (a, c)
end

function line_search(f, x, d; tol=1e-4)
    a, b = find_interval(f, x, d)
    a == b && return x
    φ = (sqrt(5.0) - 1.0) / 2.0
    c = b - φ * (b - a)
    d_pt = a + φ * (b - a)
    fc, fd = f(x + c * d), f(x + d_pt * d)
    while abs(b - a) > tol
        if fc < fd
            b, d_pt, fd = d_pt, c, fc
            c = b - φ * (b - a)
            fc = f(x + c * d)
        else
            a, c, fc = c, d_pt, fd
            d_pt = a + φ * (b - a)
            fd = f(x + d_pt * d)
        end
    end
    return x + 0.5 * (a + b) * d
end

function powell(f, x0; max_iter=50, tol=1e-4)
    n = length(x0)
    x = copy(x0)
    dirs = [[Float64(i == j) for j in 1:n] for i in 1:n]
    trajectory = [copy(x)]
    cycle_dirs = [deepcopy(dirs)]
    iters = 0
    for k in 1:max_iter
        iters = k
        x_start = copy(x)
        max_decrease, idx_replace = 0.0, 1
        for i in 1:n
            f_old = f(x)
            x = line_search(f, x, dirs[i])
            push!(trajectory, copy(x))
            decrease = f_old - f(x)
            if decrease > max_decrease
                max_decrease, idx_replace = decrease, i
            end
        end
        d_new = x - x_start
        if abs(det(hcat(dirs...))) >= 1e-6 && norm(d_new) > 1e-8
            dirs[idx_replace] = d_new ./ norm(d_new)
        end
        push!(cycle_dirs, deepcopy(dirs))
        norm(x - x_start) < tol && break
    end
    return x, trajectory, iters, cycle_dirs
end

function plot_trajectory(f, trajectory, cycle_dirs, xlim, ylim, title, filename; arrow_scale=0.15, step=1)
    n = length(trajectory[1])
    xs = range(xlim[1], xlim[2], length=100)
    ys = range(ylim[1], ylim[2], length=100)
    Z = [f([x, y]) for y in ys, x in xs]
    plt = contour(xs, ys, Z, levels=20)
    tr = trajectory
    plot!(plt, [p[1] for p in tr], [p[2] for p in tr], color=:red, linewidth=2, label="")
    scatter!(plt, [tr[1][1]], [tr[1][2]], color=:green, markersize=8, label="start")
    scatter!(plt, [tr[end][1]], [tr[end][2]], color=:blue, markersize=8, label="end")
    max_idx = length(tr) - 1
    for i in 1:step:max_idx
        cycle_idx = min(div(i - 1, n) + 1, length(cycle_dirs) - 1)
        dirs = cycle_dirs[cycle_idx]
        px, py = tr[i][1], tr[i][2]
        scale = arrow_scale * (xlim[2] - xlim[1])
        for d in dirs
            norm(d) < 1e-8 && continue
            end_x = px + d[1] * scale
            end_y = py + d[2] * scale
            plot!(plt, [px, end_x], [py, end_y], arrow=:closed, color=:blue, linewidth=1, label="")
        end
    end
    title!(plt, title)
    xlabel!(plt, "x_1")
    ylabel!(plt, "x_2")
    savefig(plt, filename)
end

x0_rastrigin = [0.3, 0.3]
x0_schwefel = [350.0, 350.0]
x0_rosenbrock = [-1.5, 2.0]

println("Функция Растригина")
x_opt, traj, iters, cycle_dirs = powell(rastrigin, x0_rastrigin)
println("  x* = ", x_opt, "\n  f(x*) = ", rastrigin(x_opt), "\n  итераций: ", iters)
plot_trajectory(rastrigin, traj, cycle_dirs, (-5, 5), (-5, 5), "Растригин", "rastrigin.png", step=1)

println("Функция Швеффеля")
x_opt, traj, iters, cycle_dirs = powell(schwefel, x0_schwefel)
println("  x* = ", x_opt, "\n  f(x*) = ", schwefel(x_opt), "\n  итераций: ", iters)
plot_trajectory(schwefel, traj, cycle_dirs, (200, 450), (200, 450), "Швеффель", "schwefel.png", step=1)

println("Функция Розенброка")
x_opt, traj, iters, cycle_dirs = powell(rosenbrock, x0_rosenbrock)
println("  x* = ", x_opt, "\n  f(x*) = ", rosenbrock(x_opt), "\n  итераций: ", iters)
plot_trajectory(rosenbrock, traj, cycle_dirs, (-2, 2), (-1, 5), "Розенброк", "rosenbrock.png", step=4)
\end{breakablelisting}

\newpage

Результат представлен на рисунках~\ref{fig:img1}--~\ref{fig:rastrigin}.

\begin{figure}[H]
  \centering
  \includegraphics[width=0.6\textwidth]{lab8/screen1.png}
  \caption{Результат}
  \label{fig:img1}
\end{figure}

\begin{figure}[H]
  \centering
  \includegraphics[width=0.6\textwidth]{lab8/screen2.png}
  \caption{Результат}
  \label{fig:schwefel}
\end{figure}

\begin{figure}[H]
  \centering
  \includegraphics[width=0.6\textwidth]{lab8/screen3.png}
  \caption{Результат}
  \label{fig:rastrigin}
\end{figure}

\end{document}
